\section{О разрешимости задачи}
	\tikzsetfigurename{FLPQ_Decidability_}

Задачи из определения~\ref{def:FLPQ} и~\ref{def:Reachability} сводятся к построению пересечения языка $\mscrL$ и языка, задаваемого путями графа, $R$.
А мы для обсуждения разрешимости задачи рассмотрим более слабую постановку задачи:

\begin{definition}[TODO: Что ты?]
    Необходимо проверить, что существует хотя бы один такой путь $\pi$ для данного графа, для данного языка $\mscrL$, что $\omega(\pi) \in \mscrL$.
\end{definition}

Эта задача сводится к проверке пустоты пересечения языка $\mscrL$ c $R$~--- регулярным языком, задаваемым графом.
От класса языка $\mscrL$ зависит её разрешимость:
\begin{itemize}
    \item Если $\mscrL$ регулярный, то получаем задачу пересечения двух регулярных языков: $\mscrL \cap R = R'$.
          $R'$~--- также регулярный язык.
          Проверка регулярного языка на пустоту~--- разрешимая проблема.
    \item Если $\mscrL$ контекстно-свободный, то получаем задачу: $\mscrL \cap R = CF$~--- контекстно-свободный.
          Проверка контекстно-свободного языка на пустоту~--- разрешимая проблема.
    \item Помимо иерархии Хомского существуют и другие классификации языков.
          Так например, класс конъюнктивных (Conj) языков~\sidecite{DBLP:journals/jalc/Okhotin01}является строгим расширением контекстно-свободных языков и все так же позволяет полиномиальный синтаксический анализ.

          Пусть $\mscrL$~--- конъюнктивный.
          При пересечении конъюнктивного и регулярного языков получается конъюнктивный ($\mscrL \cap R = Conj$), а проблема проверки Conj на пустоту не разрешима~\sidecite{DBLP:journals/tcs/Okhotin03a}.
    \item Ещё один класс языков из альтернативной иерархии, не сравнимой с Иерархией Хомского,~--- MCFG (multiple context-free grammars)~\sidecite{SEKI1991191}.
          Как его частный случай~--- TAG (tree adjoining grammar)~\sidecite{Joshi1997}.

          Если $\mscrL$ принадлежит классу MCFG, то $\mscrL \cap R$ также принадлежит MCFG.
          Проблема проверки пустоты MCFG разрешима~\sidecite{SEKI1991191}.
\end{itemize}

Существует ещё много других классификаций языков, но поиск универсальной иерархии до сих пор продолжается.

Далее, для изучения алгоритмов решения, нас будет интересовать задача $R \cap CF$.
